\chapter{Capstone Project}
\label{bab:capstone}

\section*{Tujuan Pembelajaran}

\begin{enumerate}
  \item Menjelaskan susunan sistem Exim beserta titik perluasan yang dimilikinya,
    lalu menunjuk satu fitur yang layak diperbaiki.
  \item Merumuskan dua \textit{research question} yang dapat dijawab dengan
    pengukuran, lalu menyusun metode yang menerangkan cara menjawab keduanya.
  \item Mengukur keadaan awal fitur terpilih pada tiga hal, yaitu titik sentuh,
    \textit{coupling} antar modul, dan lama satu tahap pada jalur panas.
\end{enumerate}

\section{Pendahuluan}

Empat belas bab sebelumnya memakai contoh yang sengaja dibuat kecil. Kodenya
disusun agar satu gagasan terlihat jelas, dan bebannya dirancang agar
selisihnya terukur dalam hitungan detik. Keadaan itu tidak pernah ditemui pada
sistem yang sudah dipakai orang. Sistem nyata membawa sejarah, memikul
kebergantungan yang tidak dapat dipilih sendiri, dan menolak diubah pada
bagian yang paling ingin diubah.

Bab ini menutup buku dengan satu proyek atas sistem yang benar-benar berjalan,
yaitu Exim, sebuah asisten tanya jawab ekspor impor berbasis
\textit{Retrieval-Augmented Generation} (RAG)~\cite{yohannis2026exim}. Sistem
itu terdiri atas 23 modul Python dengan 9.884 baris, ditambah 19 \textit{file}
pengujian dengan 4.868 baris. Tugasnya memilih satu fitur, memperbaiki satu
dimensi kualitasnya dengan teknik arsitektur dari bab sebelumnya, lalu
membuktikan perbaikan itu dengan angka.

Luarannya bukan kode saja, melainkan satu makalah yang disiapkan untuk jurnal
bertingkat Q3, memakai template \textit{Sakarya University Journal of Computer
and Information Sciences} (SAUCIS)~\cite{saucis2026rules}. Rancangan
pengukurannya mengikuti kaidah eksperimen perangkat
lunak~\cite{wohlin2012experimentation}, sedangkan penautan tujuan ke ukuran
mengikuti pendekatan \textit{Goal Question Metric}~\cite{basili1994gqm}. Bab
ini karena itu memuat dua hal sekaligus, yaitu cara memilih pekerjaan dan cara
melaporkannya.

\section{Sistem yang Dikerjakan}

Exim menjawab pertanyaan berbahasa Indonesia tentang kepabeanan, tarif,
larangan dan pembatasan, perizinan, serta perjanjian dagang. Jawabannya
disertai sitasi bernomor yang menunjuk \textit{file} dan halaman sumbernya. Korpusnya
berupa salinan bertanggal situs pemerintah, dan salinan itu dibuka dalam mode
baca saja.

Sistemnya tidak hanya satu. Ada tiga yang berjalan, yaitu Exim untuk ekspor
impor, satu untuk perpajakan, dan satu untuk akreditasi program studi.
Ketiganya satu garis keturunan, yaitu satu basis kode yang disalin lalu
dikembangkan sendiri-sendiri. Keadaan itu lazim disebut \textit{clone-and-own},
dan akibatnya dapat dihitung. Tabel~\ref{tab:alamat-sistem} memuat alamat
ketiganya, sedangkan Tabel~\ref{tab:tiga-sistem} membandingkan dua di antaranya.

\begin{table}[htbp]
  \centering
  \caption{Tiga sistem yang menjadi objek proyek beserta alamatnya.}
  \label{tab:alamat-sistem}
  \begin{tabularx}{\textwidth}{|l|X|X|}
    \hline
    \textbf{Sistem} & \textbf{Alamat layanan} & \textbf{Repositori} \\
    \hline
    Ekspor impor & \url{https://exim.ayolearn.stream/} & \url{https://github.com/alfa-yohannis/exim} \\
    \hline
    Perpajakan & \url{https://djp.ayolearn.stream/} & \url{https://github.com/alfa-yohannis/djp} \\
    \hline
    Akreditasi & \url{https://akreditasi.ayolearn.stream/} & \url{https://github.com/alfa-yohannis/akreditasi} \\
    \hline
  \end{tabularx}
\end{table}

Dua catatan berlaku pada tabel itu. Ketiga repositori bersifat tertutup,
sehingga akses diberikan dosen pengampu pada pertemuan pertama. Repositori
perpajakan juga belum diunggah saat bab ini ditulis, sehingga kelompok yang
memilihnya memastikan lebih dahulu bahwa repositorinya sudah tersedia.

\begin{table}[htbp]
  \centering
  \caption{Perbandingan dua turunan dari satu basis kode yang sama.}
  \label{tab:tiga-sistem}
  \begin{tabularx}{\textwidth}{|X|r|r|}
    \hline
    \textbf{Yang dibandingkan} & \textbf{Exim} & \textbf{Akreditasi} \\
    \hline
    Jumlah modul & 23 & 19 \\
    \hline
    Jumlah baris & 9.884 & 7.189 \\
    \hline
    Jumlah ketergantungan antar modul & 63 & 54 \\
    \hline
    \textit{Fan-in} modul konfigurasi & 18 & 16 \\
    \hline
    \textit{Fan-out} modul server & 12 & 10 \\
    \hline
    Modul bernama sama & \multicolumn{2}{c|}{19} \\
    \hline
    Modul yang isinya identik & \multicolumn{2}{c|}{8} \\
    \hline
    Modul yang sudah menyimpang & \multicolumn{2}{c|}{11} \\
    \hline
    Simpangan terbesar & \multicolumn{2}{c|}{\texttt{rag} 154 baris, \texttt{ingestion} 88, \texttt{server} 67} \\
    \hline
  \end{tabularx}
\end{table}

Angka itu membawa dua akibat bagi proyek ini. Pertama, satu perbaikan dapat
diukur pada lebih dari satu sistem, sehingga klaimnya tidak lagi bersandar
pada satu objek. Kedua, simpangan antar turunan itu sendiri menjadi bahan
penelitian, dan Judul 13 menggarapnya.

Jalannya data dibagi dua. Jalur pertama mengisi indeks, yaitu memindai korpus,
mengubah tiap format \textit{file} menjadi bagian teks, memotongnya menjadi potongan
yang dapat di-\textit{embed}, lalu menyimpannya pada dua tempat sekaligus.
Jalur kedua menjawab pertanyaan, yaitu mencari pada kedua \textit{storage},
menggabungkan peringkatnya, mempertajamnya dengan \textit{cross-encoder}, lalu
mengalirkan jawabannya ke peramban. Gambar~\ref{fig:exim-alur} menunjukkan
kedua jalur itu.

\begin{figure}[htbp]
  \centering
  \resizebox{\textwidth}{!}{%
  \begin{tikzpicture}[
    kotak/.style={draw, rounded corners=2pt, minimum width=2.3cm,
      minimum height=0.8cm, align=center, font=\scriptsize\bfseries},
    panah/.style={-{Stealth[length=2mm]}, thick},
    ket/.style={font=\tiny\mdseries}
  ]
    \node[kotak, fill=gray!22] (korpus) {Korpus\\
      {\tiny\mdseries baca saja}};
    \node[kotak, fill=blue!8, right=0.7cm of korpus] (ekstrak) {Extractor\\
      {\tiny\mdseries 9 format}};
    \node[kotak, fill=blue!8, right=0.7cm of ekstrak] (potong) {Chunker};
    \node[kotak, fill=blue!8, right=0.7cm of potong] (embed) {Embedder\\
      {\tiny\mdseries BGE-M3}};

    \node[kotak, fill=green!10, below=1.0cm of potong] (katalog) {Katalog\\
      {\tiny\mdseries SQLite FTS5}};
    \node[kotak, fill=green!10, below=1.0cm of embed] (vektor) {Indeks vektor\\
      {\tiny\mdseries NumPy}};

    \node[kotak, fill=praditagreen!25, below=1.0cm of katalog] (cari)
      {Hybrid retriever\\{\tiny\mdseries vektor + kata kunci}};
    \node[kotak, fill=praditagreen!25, right=0.7cm of cari] (gabung)
      {Penggabungan\\{\tiny\mdseries peringkat}};
    \node[kotak, fill=blue!8, right=0.7cm of gabung] (tajam) {Reranker\\
      {\tiny\mdseries cross-encoder}};
    \node[kotak, fill=gray!22, left=0.7cm of cari] (tanya) {Pertanyaan};
    \node[kotak, fill=gray!22, below=0.9cm of gabung] (jawab) {Jawaban\\
      {\tiny\mdseries bersitasi}};

    \draw[panah] (korpus) -- (ekstrak);
    \draw[panah] (ekstrak) -- (potong);
    \draw[panah] (potong) -- (embed);
    \draw[panah] (potong) -- (katalog);
    \draw[panah] (embed) -- (vektor);
    \draw[panah] (katalog) -- (cari);
    \draw[panah] (vektor) -- (vektor|-gabung.north);
    \draw[panah] (tanya) -- (cari);
    \draw[panah] (cari) -- (gabung);
    \draw[panah] (gabung) -- (tajam);
    \draw[panah] (tajam) |- (jawab);
  \end{tikzpicture}}
  \caption{Dua jalur pada Exim, yaitu jalur pengisian indeks di baris atas dan
    jalur menjawab pertanyaan di baris bawah.}
  \label{fig:exim-alur}
\end{figure}

Susunan modulnya dapat diukur, bukan ditaksir. Tabel~\ref{tab:exim-modul}
memuat sepuluh modul terbesar beserta \textit{fan-out} dan \textit{fan-in}.
\textit{Fan-out} adalah banyaknya modul sepaket yang diimpor satu modul,
sedangkan \textit{fan-in} adalah banyaknya modul sepaket yang mengimpornya.
Angkanya dihasilkan \texttt{ukur\_kopling.py} atas repositori yang asli.

\begin{table}[htbp]
  \centering
  \caption{Sepuluh modul terbesar Exim beserta kopling antar modulnya.}
  \label{tab:exim-modul}
  \begin{tabularx}{\textwidth}{|l|r|r|r|X|}
    \hline
    \textbf{Modul} & \textbf{Baris} & \textbf{Fan-out} & \textbf{Fan-in} & \textbf{Tanggung jawabnya} \\
    \hline
    \texttt{crawler} & 1.083 & 3 & 1 & Mengambil halaman dan lampiran situs sumber \\
    \hline
    \texttt{language\_models} & 955 & 1 & 1 & Backend Ollama dan tiga backend agen \\
    \hline
    \texttt{ingestion} & 747 & 7 & 1 & Pemindai korpus dan pipeline inkremental \\
    \hline
    \texttt{server} & 720 & 12 & 0 & Rute HTTP dan aliran \textit{Server-Sent Events} \\
    \hline
    \texttt{storage} & 691 & 3 & 4 & Katalog SQLite dan indeks vektor \\
    \hline
    \texttt{crawl\_sites} & 642 & 1 & 0 & Rencana pengambilan tiap situs \\
    \hline
    \texttt{rag} & 637 & 6 & 2 & Penyusunan \textit{prompt} dan aliran jawaban \\
    \hline
    \texttt{extractors} & 500 & 4 & 2 & Mengubah sembilan format menjadi teks \\
    \hline
    \texttt{config} & 452 & 0 & 18 & Membaca setelan menjadi objek bertipe \\
    \hline
    \texttt{retrieval} & 170 & 5 & 2 & Pencarian hibrida dan penggabungan peringkat \\
    \hline
  \end{tabularx}
\end{table}

Dua baris pada tabel itu perlu diperhatikan. Modul \texttt{config} diimpor 18
modul dari 23 yang ada, sehingga perubahan padanya merambat ke hampir seluruh
sistem. Modul \texttt{server} mengimpor 12 modul, sehingga sulit diuji
sendirian. Keduanya gejala yang dapat dihitung, dan keduanya dapat menjadi
sasaran proyek.

\section{Memilih Fitur dan Dimensi Kualitas}

Pemilihan fitur tidak dimulai dari teknik, melainkan dari gejala yang terukur.
Pertanyaannya sederhana, yaitu bagian mana yang mahal diubah, bagian mana yang
lambat, dan bagian mana yang rapuh ketika satu hal gagal. Ketiganya menjadi
tiga dimensi kualitas pada bab ini, yaitu \textit{maintainability}, performa,
dan ketahanan.

Dimensi itu menentukan ukurannya. Tabel~\ref{tab:dimensi-ukuran} memuat
ketiganya beserta ukuran yang dipakai dan cara memperolehnya. Ukuran yang
dipilih harus dapat dihitung ulang orang lain pada mesin yang berbeda, sebab
klaim yang tidak dapat diperiksa ulang akan gugur pada telaah jurnal.

Ketiga dimensi itu dipilih karena skrip pengukurannya sudah tersedia, bukan
karena hanya ada tiga. Daftar bakunya ada sembilan, yaitu model mutu produk
pada ISO/IEC 25010 edisi 2023~\cite{iso25010}. Revisi itu menambahkan
\textit{safety}, mengganti \textit{usability} menjadi \textit{interaction
capability}, mengganti \textit{portability} menjadi \textit{flexibility}, lalu
memindahkan \textit{scalability} menjadi subkarakteristik di bawah
\textit{flexibility}. Pembahasan yang lebih panjang tentang cara menalar
beberapa dimensi sekaligus ada pada laporan Software Engineering
Institute~\cite{barbacci1995quality}, sedangkan ikhtisar bidangnya ada pada
SWEBOK terbitan IEEE Computer Society~\cite{swebok2024}.

Tabel~\ref{tab:dimensi-iso} memuat kesembilannya beserta contoh yang dapat
diukur pada Exim. Kelompok boleh memilih dimensi di luar ketiga yang
disediakan, asalkan ukurannya disebut dan cara memperolehnya dapat diulang.

Daftar yang lebih luas lagi dikenal sebagai \textit{ilities}, yaitu kumpulan
sifat yang namanya banyak berakhiran sama, misalnya \textit{observability},
\textit{testability}, dan \textit{deployability}~\cite{wikipedia2026ilities}.
Daftar itu berguna untuk mencari gagasan, namun istilahnya saling tumpang
tindih dan sebagian tidak punya definisi yang disepakati. Karena itu urutannya
dijaga, yaitu gagasan dicari dari daftar luas tersebut, lalu definisi resminya
diambil dari ISO/IEC 25010 sebelum naskah ditulis.

\begin{table}[htbp]
  \centering
  \caption{Tiga dimensi kualitas beserta ukuran yang dipakai pada proyek ini.}
  \label{tab:dimensi-ukuran}
  \begin{tabularx}{\textwidth}{|l|X|X|}
    \hline
    \textbf{Dimensi} & \textbf{Ukuran} & \textbf{Cara memperolehnya} \\
    \hline
    \textit{Maintainability} & Titik sentuh, modul tersentuh, \textit{fan-in}, \textit{fan-out}, \textit{file} berubah & Pembacaan kode oleh skrip, dijalankan sebelum dan sesudah perubahan \\
    \hline
    Performa & Median, persentil ke-95, dan nilai terbesar dari lama satu tahap & Pengulangan ratusan kali pada beban yang sama \\
    \hline
    Ketahanan & Jumlah yang hilang, lama pulih, dan jumlah bagian yang ikut gagal & Mematikan satu bagian di tengah beban \\
    \hline
  \end{tabularx}
\end{table}

\begin{table}[htbp]
  \centering
  \caption{Sembilan karakteristik mutu pada ISO/IEC 25010 edisi 2023 beserta
    contoh yang dapat diukur pada Exim.}
  \label{tab:dimensi-iso}
  \begin{tabularx}{\textwidth}{|l|X|}
    \hline
    \textbf{Karakteristik} & \textbf{Contoh yang dapat diukur pada Exim} \\
    \hline
    \textit{Functional suitability} & Ketepatan sitasi, yaitu bagian jawaban yang benar-benar didukung potongan yang dikutipnya \\
    \hline
    \textit{Performance efficiency} & Latensi pencarian, hasil per detik, dan memori tiap unit pencari \\
    \hline
    \textit{Compatibility} & Berjalan di atas mesin basis data lain, atau menerima format dokumen baru \\
    \hline
    \textit{Interaction capability} & Biaya mengganti tema, menambah bahasa antarmuka, dan memenuhi syarat keterbacaan \\
    \hline
    \textit{Reliability} & Jawaban tetap selesai dan tersimpan ketika koneksi pembacanya putus \\
    \hline
    \textit{Security} & Cakupan penjagaan sesi pada seluruh rute, dan penegakan akses baca saja atas korpus \\
    \hline
    \textit{Maintainability} & Titik sentuh satu penambahan, \textit{fan-in}, dan \textit{fan-out} antar modul \\
    \hline
    \textit{Flexibility} & Biaya menambah negara, bahasa, atau backend model, termasuk \textit{scalability} ketika unit ditambah \\
    \hline
    \textit{Safety} & Jaminan korpus tidak pernah berubah, dan jawaban menyebut sumbernya pada pertanyaan regulasi \\
    \hline
  \end{tabularx}
\end{table}

\section{Dua Puluh Pilihan Judul}

Seksi ini memuat dua puluh pilihan yang seluruhnya menunjuk modul nyata pada
Exim. Seluruhnya berangkat dari gaya arsitektur yang sudah diajarkan, dan
hampir setiap bab terwakili. Tiap pilihan ditulis dengan bentuk yang sama, yaitu latar belakangnya,
pekerjaan yang dikerjakan, dua \textit{research question}, lalu cara
menjawabnya. Satu kelompok memilih satu judul, dan judul yang sama tidak boleh
dipakai dua kelompok.

Kedua puluh judul ini hanya rekomendasi. Kelompok boleh mengajukan judul lain,
asalkan judul itu memenuhi empat syarat yang sama, yaitu menunjuk modul yang
benar-benar ada pada Exim, menyebut teknik arsitektur dari bab sebelumnya,
menyebut dimensi kualitas yang diukur, dan membawa dua \textit{research
question} yang dapat dijawab dengan pengukuran. Judul usulan dibicarakan lebih
dahulu sebelum pekerjaannya dimulai, sebab judul yang ukurannya belum jelas
akan menyulitkan seksi Metode Penelitian.

Judulnya ditulis dalam bahasa Inggris, sebab naskah yang dikirim memang
berbahasa Inggris. Naskah berbahasa Indonesia yang ditulis lebih dahulu boleh
memakai terjemahannya, lalu judul Inggris itu dipakai kembali pada versi
akhirnya.

\subsection*{Judul 1: Registri Bahasa}

\textbf{Judul.} \textit{From Hard-Coded Locales to a Language Registry: Measuring the Cost of Adding a Language to a Multilingual RAG System}.

\textbf{Latar belakang.} Exim mendukung tiga bahasa jawaban, yaitu Indonesia,
Inggris, dan Vietnam. Pengetahuan tentang bahasa itu tersebar, yaitu 26 titik
sentuh pada \berkas{app/rag.py} dan satu titik pada \berkas{app/server.py}
berupa tipe \texttt{Literal["id", "en", "vi"]}.

\textbf{Yang dikerjakan.} Pengetahuan bahasa dipusatkan ke satu registri,
sehingga menambah bahasa cukup menambah satu entri, bukan menyunting beberapa
\textit{file}.

\begin{enumerate}
  \item Berapa titik sentuh dan modul yang harus disunting untuk menambah satu
    bahasa, sebelum dan sesudah registri dipakai?
  \item Apakah pemusatan itu menghilangkan risiko terlewat, diukur dari titik
    sentuh yang masih tersisa di luar registri?
\end{enumerate}

\textbf{Cara menjawab.} RQ1 dijawab dengan menjalankan
\texttt{ukur\_titik\_sentuh.py} pada kedua keadaan lalu membandingkan
angkanya. RQ2 dijawab dengan menambah satu bahasa keempat sungguhan, lalu
mencatat \textit{file} yang berubah dan hasil seluruh pengujian.

\subsection*{Judul 2: Manifes Situs}

\textbf{Judul.} \textit{Adding a Country to a Regulatory RAG Corpus: A Plugin Manifest for Crawl Sites and Its Effect on Change Containment}.

\textbf{Latar belakang.} Korpus Exim memuat sepuluh situs dari dua negara,
yaitu sembilan situs Indonesia dan satu situs Vietnam. Menambah satu negara
hari ini menyentuh delapan modul dan 72 baris.

\textbf{Yang dikerjakan.} Rencana pengambilan tiap situs dideklarasikan lewat
manifes, sehingga penambahan negara tidak lagi menyunting modul yang mengenal
isi negara tertentu.

\begin{enumerate}
  \item Berapa modul yang ikut berubah ketika satu negara ditambahkan, sebelum
    dan sesudah manifes dipakai?
  \item Apakah satu situs yang gagal diambil menghentikan pengambilan situs
    lain?
\end{enumerate}

\textbf{Cara menjawab.} RQ1 dijawab dengan \texttt{ukur\_titik\_sentuh.py}
memakai nama negara sebagai penanda. RQ2 dijawab dengan menjalankan
pengambilan yang satu situsnya sengaja dibuat gagal, lalu mencatat situs lain
yang tetap selesai.

\subsection*{Judul 3: Backend Model sebagai Plugin}

\textbf{Judul.} \textit{Plugin-Based Model Backends for a Local RAG System: Measuring the Cost of Adding and Replacing Language Models}.

\textbf{Latar belakang.} Modul \berkas{app/language\_models.py} berisi 955
baris dengan satu kelas abstrak dan beberapa turunan, yaitu backend Ollama dan
backend agen baris perintah. Penambahan backend masih menyentuh beberapa
\textit{file}.

\textbf{Yang dikerjakan.} Backend didaftarkan lewat registri beserta
manifesnya, lalu dimuat hanya ketika dipakai.

\begin{enumerate}
  \item Berapa \textit{file} dan baris yang berubah untuk menambah satu backend baru,
    sebelum dan sesudah registri dipakai?
  \item Apakah pemuatan malas menurunkan \textit{startup
    time}\index{startup time@\textit{startup time}} tanpa menaikkan
    latensi pemakaian pertama?
\end{enumerate}

\textbf{Cara menjawab.} RQ1 dijawab dengan menambah satu backend tiruan lalu
menghitung \textit{file} yang berubah. RQ2 dijawab dengan mengukur lama penyalaan dan
lama panggilan pertama, masing-masing seratus putaran.

\subsection*{Judul 4: Kopling Modul Konfigurasi}

\textbf{Judul.} \textit{Reducing Coupling of a Configuration Hub: Fan-In Before and After Interface Segregation}.

\textbf{Latar belakang.} Modul \berkas{app/config.py} diimpor 18 modul dari 23
yang ada. Setiap perubahan padanya karena itu berpeluang merambat ke hampir
seluruh sistem.

\textbf{Yang dikerjakan.} Setelan dipecah menurut pemakainya, sehingga tiap
modul hanya bergantung pada bagian setelan yang benar-benar dipakainya.

\begin{enumerate}
  \item Berapa \textit{fan-in} modul konfigurasi sebelum dan sesudah
    pemecahan?
  \item Berapa modul yang ikut dimuat ketika satu modul diuji sendirian,
    sebelum dan sesudah pemecahan?
\end{enumerate}

\textbf{Cara menjawab.} RQ1 dijawab dengan \texttt{ukur\_kopling.py}. RQ2
dijawab dengan mencatat isi \texttt{sys.modules} sesudah satu modul diimpor,
lalu membandingkan jumlahnya pada kedua keadaan.

\subsection*{Judul 5: Memecah Modul HTTP}

\textbf{Judul.} \textit{Splitting an HTTP God-Module: Measuring Fan-Out and Change Containment in a RAG Server}.

\textbf{Latar belakang.} Modul \berkas{app/server.py} berisi 720 baris dan
mengimpor 12 modul, dan itu \textit{fan-out} terbesar pada sistem ini.

\textbf{Yang dikerjakan.} Rute dipisah menurut kemampuannya, lalu tiap bagian
hanya memanggil modul yang dibutuhkannya.

\begin{enumerate}
  \item Berapa \textit{fan-out} tiap bagian sesudah rute dipisah, dibandingkan
    satu modul tunggal sebelumnya?
  \item Berapa \textit{file} yang berubah ketika satu rute baru ditambahkan, sebelum
    dan sesudah pemisahan?
\end{enumerate}

\textbf{Cara menjawab.} RQ1 dijawab dengan \texttt{ukur\_kopling.py} atas
paket yang sudah dipecah. RQ2 dijawab dengan menambah satu rute sungguhan lalu
menghitung \textit{file} yang disentuh.

\subsection*{Judul 6: Registry Extractor}

\textbf{Judul.} \textit{The Extractor Registry as a Microkernel\index{microkernel}: The
Cost of Adding a New Document Format}.

\textbf{Latar belakang.} Modul \berkas{app/extractors.py} memuat sembilan
extractor di bawah satu registri. Penambahan format baru masih menyunting
\textit{file} yang sama dengan seluruh extractor lain.

\textbf{Yang dikerjakan.} Tiap extractor dipindahkan menjadi plugin dengan
manifesnya sendiri, lalu dimuat ketika formatnya ditemui.

\begin{enumerate}
  \item Berapa \textit{file} yang berubah untuk menambah satu format baru, sebelum
    dan sesudah plugin dipakai?
  \item Apakah satu extractor yang gagal dimuat menghentikan pemuatan extractor
    lain?
\end{enumerate}

\textbf{Cara menjawab.} RQ1 dijawab dengan menambah satu extractor sungguhan.
RQ2 dijawab dengan menyisipkan satu plugin yang sengaja gagal, lalu mencatat
extractor lain yang tetap tersedia.

\subsection*{Judul 7: Adapter Indeks Vektor}

\textbf{Judul.} \textit{Ports and Adapters for the Vector Index: Replacing Exhaustive Search Without Touching the Core}.

\textbf{Latar belakang.} Pencarian vektor Exim memakai perkalian titik
menyeluruh, dan dokumen arsitekturnya menyebut \texttt{VectorIndex.search}
sebagai satu-satunya metode yang perlu diganti bila korpus tumbuh.

\textbf{Yang dikerjakan.} Klaim itu diuji, yaitu adapter kedua dipasang di
belakang \textit{port}\index{port@\textit{port}} yang sama, lalu keduanya diukur pada
korpus tiruan
dengan ukuran bertingkat.

\begin{enumerate}
  \item Berapa latensi pencarian kedua adapter pada ukuran korpus yang
    berbeda?
  \item Berapa \textit{file} di luar adapter yang ikut berubah ketika adapter ditukar?
\end{enumerate}

\textbf{Cara menjawab.} RQ1 dijawab dengan seratus putaran pencarian pada tiap
ukuran korpus, lalu melaporkan median\index{median}, persentil
ke-95\index{persentil ke-95}, dan nilai terbesar.
RQ2 dijawab dengan \texttt{git diff} yang dikelompokkan menurut \textit{file}-nya.

\subsection*{Judul 8: Urutan Tahap Pencarian}

\textbf{Judul.} \textit{Stage Reordering in a Hybrid Retrieval Pipeline: Where Should the Folder Filter Run?}

\textbf{Latar belakang.} Pencarian Exim berjalan empat tahap, yaitu pencarian
vektor, pencarian kata kunci, penggabungan peringkat, dan \textit{reranking}.
Penyaring folder dijalankan sebelum pencarian vektor.

\textbf{Yang dikerjakan.} Urutan tahapnya diubah, lalu akibatnya pada lama
pencarian dan pada hasil yang keluar dibandingkan.

\begin{enumerate}
  \item Berapa lama tiap tahap pada kedua urutan?
  \item Apakah hasil sepuluh teratas berubah ketika urutannya diubah?
\end{enumerate}

\textbf{Cara menjawab.} RQ1 dijawab dengan mengukur tiap tahap secara
terpisah, seratus putaran. RQ2 dijawab dengan membandingkan daftar
\texttt{chunk\_id} kedua urutan memakai pertanyaan yang sama.

\subsection*{Judul 9: Indeks yang Direplikasi Antar Unit Pencari}

\textbf{Judul.} \textit{A Replicated In-Memory Index Across Query Units:
Retrieval Throughput and Memory Cost Under Simultaneous Users}.

\textbf{Latar belakang.} Indeks vektor Exim berupa matriks di memori proses,
dan dokumen arsitekturnya menaksir ukurannya sekitar 82 MB untuk 20.000
potongan. Satu proses menjawab satu pencarian pada satu waktu, sehingga
pemakai yang datang serentak mengantre.

\textbf{Yang dikerjakan.} Beberapa unit pencari dijalankan, masing-masing
memegang copy indeksnya sendiri, mengikuti susunan pada
Bab~\ref{bab:ruang}. Beban Exim adalah keadaan yang menguntungkan bagi grid
yang direplikasi, sebab isinya banyak dibaca dan hanya berubah ketika korpus
diindeks ulang.

\begin{enumerate}
  \item Berapa pencarian per detik yang dilayani pada satu sampai enam unit,
    di bawah pemakai yang datang serentak?
  \item Berapa memori yang dipakai tiap unit tambahan, dan pada jumlah unit
    berapa memorinya habis?
\end{enumerate}

\textbf{Cara menjawab.} RQ1 dijawab dengan mengirim beban serentak ke
sejumlah unit, lima putaran tiap jumlah unit, lalu melaporkan median beserta
persentil ke-95. RQ2 dijawab dengan mencatat pemakaian memori tiap proses.
Perlu disebutkan pada Pembahasan bahwa lama menulis jawaban tetap ditentukan
satu GPU, sehingga yang membaik hanya jalur pencariannya.

\subsection*{Judul 10: Pipeline Indexing Inkremental}

\textbf{Judul.} \textit{Pipe-and-Filter Ingestion: Throughput and First-Output Latency of an Incremental Indexer}.

\textbf{Latar belakang.} Modul \berkas{app/ingestion.py} berisi 747 baris dan
mengimpor tujuh modul. Jalan pertama atas korpus memakan puluhan menit,
sedangkan jalan kedua hanya membaca \textit{file} yang berubah.

\textbf{Yang dikerjakan.} Tahapnya dijalankan berdampingan, lalu hasilnya
dibandingkan dengan tahap yang berjalan berurutan.

\begin{enumerate}
  \item Berapa \textit{file} per detik yang selesai diindeks pada kedua susunan?
  \item Berapa lama sampai potongan pertama tersimpan pada kedua susunan?
\end{enumerate}

\textbf{Cara menjawab.} Keduanya dijawab dengan korpus tiruan berisi ratusan
\textit{file} kecil, lima putaran, lalu melaporkan mediannya beserta ragamnya.

\subsection*{Judul 11: Ketahanan Aliran Jawaban}

\textbf{Judul.} \textit{Event-Driven Answer Streaming: Measuring Resilience of Resumable Jobs Under Connection Loss}.

\textbf{Latar belakang.} Jawaban Exim dibuat pada pekerjaan latar belakang
yang mencatat tiap event dengan nomor urut, sehingga pembaca yang terputus
dapat menyambung dari nomor terakhir yang diterimanya.

\textbf{Yang dikerjakan.} Ketahanan itu diuji, yaitu koneksi diputus pada
berbagai titik, lalu kelengkapan jawabannya diperiksa.

\begin{enumerate}
  \item Berapa event yang hilang ketika koneksi diputus di tengah jawaban?
  \item Berapa lama pembaca kembali menerima event sesudah menyambung
    ulang?
\end{enumerate}

\textbf{Cara menjawab.} Keduanya dijawab dengan memutus koneksi pada sepuluh
titik berbeda, masing-masing sepuluh putaran, lalu membandingkan nomor
event yang diterima dengan nomor yang tercatat.

\subsection*{Judul 12: Dua \textit{Storage}}

\textbf{Judul.} \textit{Building a Read Model for Hybrid Retrieval: Query
Latency Against Write Cost and the Stale-Read Window}.

\textbf{Latar belakang.} Baris potongan disimpan ke SQLite per dokumen,
sedangkan arsip vektor disimpan setiap 25 dokumen. Proses yang terhenti di
antara keduanya meninggalkan dokumen yang tercatat terindeks tetapi vektornya
belum ada.

\textbf{Yang dikerjakan.} Sisi baca diberi model tersendiri yang memang
disiapkan untuk kueri, misalnya himpunan \texttt{chunk\_id} per folder yang
dihitung di muka, mengikuti pemisahan sisi baca dan sisi tulis pada
Bab~\ref{bab:cqrs}. Harganya lalu diukur, yaitu tambahan pekerjaan di sisi
tulis dan lebar jendela ketika kedua sisi belum sejalan.

\begin{enumerate}
  \item Berapa latensi pencarian dengan dan tanpa model baca itu?
  \item Berapa tambahan waktu di sisi tulis, dan berapa lebar jendela
    ketidaksesuaian pada beberapa ukuran kelompok \textit{storage}?
\end{enumerate}

\textbf{Cara menjawab.} RQ1 dijawab dengan menghentikan proses pada titik
acak, lalu menghitung dokumen yang kehilangan vektor. RQ2 dijawab dengan
mengukur lama pemeriksaan awal pada tiap ukuran kelompok.

\subsection*{Judul 13: Inti Bersama bagi Tiga Turunan}

\textbf{Judul.} \textit{From Clone-and-Own to a Shared Core: Measuring
Divergence and Merge Cost Across Three Deployments of One RAG Codebase}.

\textbf{Latar belakang.} Ketiga sistem berasal dari satu basis kode yang
disalin lalu dikembangkan terpisah. Dari 19 modul bernama sama antara Exim dan
Akreditasi, delapan masih identik baris per baris sedangkan sebelas sudah
menyimpang, dengan simpangan terbesar 154 baris pada \texttt{rag}. Perbaikan
pada satu turunan karena itu tidak sampai ke turunan lain.

\textbf{Yang dikerjakan.} Modul yang masih identik diangkat menjadi inti
bersama, sedangkan modul yang menyimpang diperiksa, yaitu apakah simpangannya
memang kebutuhan yang berbeda atau sekadar perbaikan yang belum menyeberang.
Bagian yang memang berbeda dinyatakan sebagai titik perluasan, mengikuti
Bab~\ref{bab:plugin}.

\begin{enumerate}
  \item Berapa baris yang dapat disatukan menjadi inti bersama, dan berapa
    yang memang harus tetap berbeda?
  \item Berapa biaya menerapkan satu perbaikan ke ketiga sistem, sebelum dan
    sesudah inti bersama dipakai?
\end{enumerate}

\textbf{Cara menjawab.} RQ1 dijawab dengan membandingkan isi tiap modul
bernama sama, lalu menggolongkan tiap blok yang berbeda menjadi kebutuhan yang
berbeda atau perbaikan yang tertinggal. RQ2 dijawab dengan menerapkan satu
perbaikan sungguhan, misalnya satu perbaikan keamanan, lalu menghitung \textit{file}
dan baris yang disentuh pada kedua susunan.

\textbf{\textit{Trade-off} yang wajib dilaporkan.}\index{trade-off@\textit{trade-off}}
Inti bersama mengurangi
penggandaan, namun mengikat ketiga sistem pada satu jadwal rilis. Perubahan
pada inti menuntut pengujian ketiganya, dan satu kebutuhan khas satu sistem
menjadi lebih mahal dipenuhi. Kedua sisi itu dilaporkan berangka, yaitu baris
yang hemat dan jumlah pengujian yang ikut berjalan.

\subsection*{Judul 14: Arah Ketergantungan Antar Layer}

\textbf{Judul.} \textit{Does the Dependency Flow Really Go One Way? Detecting
and Removing Layering Violations in a RAG Codebase}.

\textbf{Latar belakang.} Dokumen arsitektur Exim menyatakan bahwa
ketergantungan mengalir satu arah, yaitu \texttt{storage} tidak mengenal HTTP
dan \texttt{server} tidak menyentuh SQLite langsung. Pernyataan itu dapat
dibuktikan salah, sehingga layak diuji.

\textbf{Yang dikerjakan.} Urutan layer ditetapkan lebih dahulu, lalu seluruh
impor diperiksa terhadap urutan itu, dan pelanggaran yang ditemukan
dihilangkan.

\begin{enumerate}
  \item Berapa pelanggaran arah ketergantungan yang ada pada keadaan awal?
  \item Berapa \textit{file} yang berubah untuk menghilangkannya, dan apakah seluruh
    pengujian tetap lulus?
\end{enumerate}

\textbf{Cara menjawab.} RQ1 dijawab dengan memperluas
\texttt{ukur\_kopling.py} agar membandingkan nomor layer pengimpor dengan
nomor layer yang diimpor. RQ2 dijawab dengan \texttt{git diff} beserta hasil
seluruh pengujian.

\subsection*{Judul 15: Memisahkan Indexing dari Pelayanan}

\textbf{Judul.} \textit{Splitting Ingestion from Query Serving: Measuring
Deployment Independence and Startup Time in a Service-Based Split}.

\textbf{Latar belakang.} Pengisian indeks dan pelayanan pertanyaan berjalan di
dalam satu proses yang sama, padahal keduanya punya beban dan jadwal yang
berbeda.

\textbf{Yang dikerjakan.} Keduanya dipisah menjadi dua layanan yang berbagi
\textit{storage}, mengikuti susunan \textit{service-based} pada Bab~\ref{bab:layanan}.

\begin{enumerate}
  \item Berapa \textit{startup time} layanan pelayanan sesudah bagian indexing
    dikeluarkan?
  \item Berapa bagian yang ikut berhenti ketika layanan indexing dimatikan di
    tengah beban?
\end{enumerate}

\textbf{Cara menjawab.} RQ1 dijawab dengan seratus putaran penyalaan pada
kedua susunan. RQ2 dijawab dengan mematikan layanan indexing lalu mencatat
pertanyaan yang tetap terjawab.

\subsection*{Judul 16: Crawler sebagai Unit Tersendiri}

\textbf{Judul.} \textit{Extracting the Crawler as an Independently Deployable
Service: Measuring Build Time, Test Time, and Failure Isolation}.

\textbf{Latar belakang.} Modul \berkas{app/crawler.py} berisi 1.083 baris,
modul terbesar pada Exim, dan hanya dipakai ketika korpus diperbarui.

\textbf{Yang dikerjakan.} Crawler dikeluarkan menjadi unit yang dapat
diterbitkan sendiri, mengikuti pertimbangan pada Bab~\ref{bab:mikroservis}.

\begin{enumerate}
  \item Berapa lama pengujian berjalan sesudah crawler dikeluarkan,
    dibandingkan sebelumnya?
  \item Berapa kemampuan yang ikut gagal ketika crawler gagal berjalan, pada
    kedua susunan?
\end{enumerate}

\textbf{Cara menjawab.} RQ1 dijawab dengan menjalankan seluruh pengujian lima
kali pada kedua susunan. RQ2 dijawab dengan membuat crawler gagal dengan
sengaja, lalu mencatat kemampuan yang tetap tersedia.

\subsection*{Judul 17: Mengganti Mesin Basis Data}

\textbf{Judul.} \textit{Swapping the Storage Engine Behind a Port: How Leaky
Is an FTS5-Shaped Abstraction?}

\textbf{Latar belakang.} Seluruh \textit{storage} Exim memakai SQLite. Modul
\berkas{app/storage.py} memuat 17 baris yang bergantung pada fitur khas mesin
itu, yaitu sintaks FTS5, fungsi \texttt{bm25()}, dan beberapa \texttt{PRAGMA}.
Kelas \texttt{KeywordQueryBuilder} bahkan seluruhnya menyusun ekspresi
\texttt{MATCH} milik FTS5, sehingga pencarian kata kuncinya tidak dapat
dipindahkan begitu saja.

\textbf{Yang dikerjakan.} Pencarian kata kunci dinyatakan sebagai satu
\textit{port}, lalu adapter kedua dibuat pada mesin basis data lain, misalnya
PostgreSQL dengan \texttt{tsvector}, mengikuti Bab~\ref{bab:port}.

\begin{enumerate}
  \item Berapa baris dan \textit{file} yang masih bergantung pada fitur khas satu
    mesin sesudah \textit{port} dipasang?
  \item Berapa latensi kueri kedua adapter pada ukuran korpus yang sama, dan
    berapa besar irisan sepuluh hasil teratasnya?
\end{enumerate}

\textbf{Cara menjawab.} RQ1 dijawab dengan \texttt{ukur\_titik\_sentuh.py}
memakai penanda \texttt{sqlite3}, \texttt{fts5}, \texttt{bm25}, dan
\texttt{PRAGMA}. RQ2 dijawab dengan seratus kueri yang sama pada korpus
tiruan, lalu melaporkan median beserta persentil ke-95, ditambah irisan daftar
\texttt{chunk\_id} kedua adapter.

\subsection*{Judul 18: Pergantian Tema Tampilan}

\textbf{Judul.} \textit{Theme Tokens Across Two Stylesheets: Measuring the
Cost of Swapping a Visual Theme}.

\textbf{Latar belakang.} Tampilan Exim memakai dua lembar gaya, yaitu
\berkas{web/styles.css} dengan 1.878 baris dan \berkas{web/landing.css} dengan
1.139 baris. Keduanya sudah memakai token warna, namun warna heksadesimal masih
muncul pada 167 baris yang tersebar di enam \textit{file}, dan tokennya dideklarasikan
dua kali pada dua \textit{file} terpisah.

\textbf{Yang dikerjakan.} Token tema disatukan menjadi satu sumber yang
dipakai kedua halaman, lalu warna yang masih dipasang keras dipindahkan ke
token itu, mengikuti pemisahan tampilan dari isi pada Bab~\ref{bab:mv}.

\begin{enumerate}
  \item Berapa warna yang masih dipasang keras di luar token, sebelum dan
    sesudah penyatuan?
  \item Berapa \textit{file} dan baris yang harus disunting untuk menambah satu tema
    baru, sebelum dan sesudahnya?
\end{enumerate}

\textbf{Cara menjawab.} RQ1 dijawab dengan \texttt{ukur\_titik\_sentuh.py}
memakai penanda warna heksadesimal pada direktori \berkas{web}. RQ2 dijawab
dengan menambah satu tema gelap sungguhan, lalu menghitung \textit{file} dan baris
yang disentuh.

\subsection*{Judul 19: Penjagaan Sesi pada Rute}

\textbf{Judul.} \textit{Authentication as a Layer: Measuring Guard Coverage
and the Cost of Adding a Protected Route}.

\textbf{Latar belakang.} Exim memeriksa sesi lewat satu \textit{middleware}
HTTP, yaitu \texttt{\_require\_session} pada \berkas{app/server.py}.
Pemeriksaan itu berlaku bagi seluruh rute, sehingga rute yang memang harus
terbuka perlu dikecualikan satu per satu.

\textbf{Yang dikerjakan.} Penjagaan itu dipindahkan menjadi kemampuan yang
dinyatakan tiap rute, bukan daftar perkecualian, mengikuti pembagian tanggung
jawab antar layer pada Bab~\ref{bab:lapisan}.

\begin{enumerate}
  \item Berapa rute yang benar-benar terjaga dan berapa yang terbuka, pada
    kedua susunan?
  \item Berapa \textit{file} yang berubah untuk menambah satu rute terjaga, sebelum
    dan sesudahnya?
\end{enumerate}

\textbf{Cara menjawab.} RQ1 dijawab dengan memanggil seluruh rute tanpa
\textit{cookie} sesi, lalu mencatat yang membalas 401 dan yang tidak. RQ2
dijawab dengan menambah satu rute terjaga sungguhan.

\subsection*{Judul 20: Penambahan User dan Role}

\textbf{Judul.} \textit{From Two Hard-Coded Roles to a Permission Registry:
Measuring the Cost of Adding a User and a Role}.

\textbf{Latar belakang.} Exim mengenal dua \textit{role} saja, yaitu
\texttt{admin} dan \texttt{guest}, keduanya ditulis sebagai konstanta pada
\berkas{app/authentication.py}. Akunnya satu untuk tiap \textit{role}, diambil
dari setelan, dan izinnya berupa satu boolean \texttt{is\_admin} yang muncul
pada empat baris di dua \textit{file}. Menambah pemakai atau \textit{role} karena itu
menuntut penyuntingan kode, bukan penambahan data.

\textbf{Yang dikerjakan.} Pemakai dan \textit{role} dipindahkan menjadi data,
lalu izin dinyatakan sebagai daftar kemampuan pada satu registri, mengikuti
pemisahan \textit{core}\index{core@\textit{core}} dari pengisinya pada
Bab~\ref{bab:port} dan
Bab~\ref{bab:plugin}.

\begin{enumerate}
  \item Berapa \textit{file} dan baris yang berubah untuk menambah satu pemakai dan
    satu \textit{role} baru, sebelum dan sesudah registri dipakai?
  \item Berapa tempat pemeriksaan izin yang masih memakai boolean
    \texttt{is\_admin} sesudah perubahan?
\end{enumerate}

\textbf{Cara menjawab.} RQ1 dijawab dengan menambah satu \textit{role}
sungguhan, misalnya peninjau yang boleh membaca tetapi tidak boleh mengubah
indeks, lalu menghitung \textit{file} yang disentuh. RQ2 dijawab dengan
\texttt{ukur\_titik\_sentuh.py} memakai \texttt{is\_admin} sebagai penanda.

\section{Pola Rancangan sebagai Alat}

Gaya arsitektur menentukan bentuk besar sebuah sistem, sedangkan pola
rancangan menentukan cara memasangnya pada tingkat kelas dan
antarmuka~\cite{gamma1994patterns}. Keduanya saling melengkapi, dan sebagian
pola bahkan dirumuskan khusus untuk tingkat arsitektur~\cite{buschmann1996posa}.
Kedua puluh judul pada seksi sebelumnya seluruhnya diwujudkan dengan pola,
meski namanya belum disebut pada judulnya.

Kedudukan pola pada makalah perlu dijaga. Pola disebut pada seksi Desain
Solusi sebagai cara mewujudkan rancangan, bukan pada seksi Pendahuluan sebagai
sumbangan. Makalah yang sumbangannya hanya berupa penerapan satu pola akan
ditolak, sebab pola itu sudah dikenal luas sejak lama. Yang menjadi sumbangan
tetap akibat terukurnya pada dimensi kualitas yang dipilih.

Pola juga membawa harganya sendiri. Setiap \textit{layer}\index{layer@\textit{layer}}
tak langsung menambah satu
panggilan, dan tambahan itu terukur. Makalah yang melaporkan bahwa polanya
memperbaiki \textit{maintainability} sekaligus menaikkan latensi sekian
mikrodetik justru lebih kuat, sebab pembaca melihat kedua sisinya.
Tabel~\ref{tab:pola} memuat pola yang lazim dipakai tiap kelompok judul.

\begin{table}[htbp]
  \centering
  \caption{Pola rancangan yang mewujudkan tiap kelompok judul.}
  \label{tab:pola}
  \begin{tabularx}{\textwidth}{|l|X|X|}
    \hline
    \textbf{Judul} & \textbf{Pola yang dipakai} & \textbf{Harga yang ditimbulkannya} \\
    \hline
    1, 2, 3, 6, 13 & \textit{Registry} dan \textit{Factory} untuk pendaftaran, \textit{Strategy} untuk perilaku yang berbeda tiap entri & Satu \textit{layer} pencarian entri sebelum pekerjaan sebenarnya dijalankan \\
    \hline
    4, 5 & \textit{Facade} per pemakai dan \textit{Parameter Object} untuk setelan yang dipecah & Jumlah \textit{file} bertambah, sehingga pembaca baru perlu menelusuri lebih banyak \textit{file} \\
    \hline
    7 & \textit{Adapter} di belakang satu \textit{port}, dengan \textit{Strategy} memilih implementasinya & Satu panggilan tak langsung pada tiap pencarian \\
    \hline
    8, 10 & \textit{Pipes and Filters} dengan tiap tahap berdiri sendiri & Penyerahan data antar tahap, yang pada aliran besar menjadi biaya nyata \\
    \hline
    9 & \textit{Proxy} yang menyimpan hasil, dibungkus \textit{Decorator} & Jawaban dapat memakai hasil lama, yaitu \textit{stale read} \\
    \hline
    11, 18 & \textit{Observer} atas catatan event bernomor urut & Catatan event harus disimpan, sehingga memakan memori \\
    \hline
    12 & \textit{Repository} terpisah untuk sisi baca dan sisi tulis & Kedua \textit{storage} dapat tidak sejalan, sehingga perlu rekonsiliasi \\
    \hline
    14, 15, 16, 17, 19, 20 & \textit{Facade} pada batas layanan, dan \textit{Template Method} pada langkah yang berulang & Batas baru menambah satu \textit{hop}, dan \textit{hop} itu terukur \\
    \hline
  \end{tabularx}
\end{table}

Satu catatan praktis. Nama pola ditulis apa adanya, bukan diterjemahkan, sebab
penelaah mengenalinya dari nama aslinya. Penyebutannya juga disertai alasan,
misalnya \textit{Adapter}\index{adapter@\textit{adapter}} dipilih karena antarmuka
lamanya tidak boleh berubah,
bukan sekadar karena polanya terkenal.

\section{Merumuskan Research Question}

Pertanyaan penelitian yang baik menyebut tiga hal sekaligus, yaitu apa yang
diubah, apa yang diukur, dan pada keadaan mana. Pertanyaan yang hanya menyebut
satu di antaranya tidak dapat dijawab dengan angka. Kalimat seperti "apakah
arsitektur ini lebih baik" tidak memenuhi syarat, sebab kata lebih baik tidak
menyebut ukurannya.

Tiap kelompok wajib menyusun sedikitnya dua pertanyaan. Keduanya tidak boleh
menanyakan hal yang sama dengan kata yang berbeda. Tabel~\ref{tab:bentuk-rq}
memuat bentuk yang dipakai beserta contoh ukurannya.

\begin{table}[htbp]
  \centering
  \caption{Bentuk \textit{research question} beserta ukuran dan buktinya.}
  \label{tab:bentuk-rq}
  \begin{tabularx}{\textwidth}{|l|X|X|}
    \hline
    \textbf{Bentuk} & \textbf{Contoh pertanyaan} & \textbf{Bukti yang menjawabnya} \\
    \hline
    Biaya perubahan & Berapa \textit{file} dan baris yang berubah untuk menambah satu bahasa? & Keluaran skrip pembaca kode pada dua keadaan \\
    \hline
    Lama pekerjaan & Berapa median dan persentil ke-95 lama satu tahap? & Ratusan putaran pada beban yang sama \\
    \hline
    Ketahanan & Berapa bagian yang ikut gagal ketika satu bagian gagal? & Kegagalan yang sengaja disisipkan \\
    \hline
    Ketepatan & Apakah hasil sepuluh teratas berubah sesudah susunannya diubah? & Perbandingan daftar keluaran pada pertanyaan yang sama \\
    \hline
  \end{tabularx}
\end{table}

\section{Metode Penelitian}

Metode Penelitian pada makalah wajib menerangkan cara tiap \textit{research
question} dijawab, satu per satu. Penjelasan itu bukan kalimat umum tentang
pendekatan kuantitatif, melainkan urutan langkah yang dapat diulang orang
lain. Pembaca harus dapat menjalankan ulang seluruhnya hanya dengan membaca
seksi itu.

Rancangannya memakai bentuk sebelum dan sesudah pada sistem yang sama.
Keadaan awal diukur lebih dahulu, perubahan arsitekturnya diterapkan, lalu
keadaan akhir diukur dengan cara yang persis sama.
Gambar~\ref{fig:rancangan-ukur} menunjukkan urutannya.

\begin{figure}[htbp]
  \centering
  \resizebox{\textwidth}{!}{%
  \begin{tikzpicture}[
    kotak/.style={draw, rounded corners=2pt, minimum width=2.6cm,
      minimum height=0.9cm, align=center, font=\scriptsize\bfseries},
    panah/.style={-{Stealth[length=2mm]}, thick},
    ket/.style={font=\tiny\mdseries}
  ]
    \node[kotak, fill=gray!22] (awal) {Keadaan awal\\
      {\tiny\mdseries branch induk}};
    \node[kotak, fill=blue!8, right=1.1cm of awal] (ukur1) {Pengukuran 1\\
      {\tiny\mdseries 100 putaran}};
    \node[kotak, fill=praditagreen!25, right=1.1cm of ukur1] (ubah)
      {Perubahan\\{\tiny\mdseries teknik arsitektur}};
    \node[kotak, fill=blue!8, right=1.1cm of ubah] (ukur2) {Pengukuran 2\\
      {\tiny\mdseries skrip yang sama}};
    \node[kotak, fill=green!10, below=1.0cm of ubah] (banding)
      {Perbandingan\\{\tiny\mdseries selisih lawan ragam}};

    \draw[panah] (awal) -- (ukur1);
    \draw[panah] (ukur1) -- (ubah);
    \draw[panah] (ubah) -- (ukur2);
    \draw[panah] (ukur1) |- (banding);
    \draw[panah] (ukur2) |- (banding);
  \end{tikzpicture}}
  \caption{Rancangan sebelum dan sesudah pada sistem yang sama, dengan skrip
    pengukuran yang tidak diubah di antara keduanya.}
  \label{fig:rancangan-ukur}
\end{figure}

Tiga aturan berlaku pada pengukurannya. Pertama, satu kali pengukuran tidak
cukup, sehingga tiap keadaan diulang sedikitnya seratus kali. Kedua, yang
dilaporkan bukan rata-rata, melainkan median, persentil ke-95, dan nilai
terbesar, sebab gejala antrean hanya terlihat pada nilai terbesar. Ketiga,
selisih yang lebih kecil daripada ragam antar putaran belum dapat disebut
ada~\cite{wohlin2012experimentation}.

Pengukuran yang membaca kode mengikuti aturan yang berbeda. Angkanya tidak
beragam, sehingga tidak perlu diulang. Yang wajib disebutkan justru definisinya,
misalnya apakah komentar ikut dihitung sebagai titik sentuh, sebab definisi
yang berbeda menghasilkan angka yang berbeda.

Tiga skrip disediakan sebagai titik berangkat, dan ketiganya hanya memakai
pustaka standar Python. Listing~\ref{lst:keadaan-awal} menunjukkan cara
menjalankannya beserta keluaran nyatanya pada repositori Exim.

\begin{lstlisting}[language=bash, caption={Mengukur keadaan awal Exim, tersedia di \berkas{sources/chapter15/README.md}}, label={lst:keadaan-awal}]
git clone https://github.com/alfa-yohannis/exim.git
cd sources/chapter15

python ukur_kopling.py exim/app
# Keluarannya:
# Jumlah modul        : 23
# Jumlah baris        : 9884
# Fan-in terbesar     : config, 18 modul mengimpornya
# Fan-out terbesar    : server, mengimpor 12 modul

python ukur_titik_sentuh.py exim/app '"id"' '"en"' '"vi"'
# Keluarannya:
# File tersentuh      : 2
# Baris tersentuh     : 27

python ukur_penggabungan.py
# Keluarannya:
#  Kandidat masuk  Kandidat keluar   Median us    P95 us   Terbesar us
#             100               85       19.17     21.49         45.87
#             400              340       73.45     89.22        129.09
\end{lstlisting}

Keluaran terakhir itu sekaligus menjadi contoh \textit{benchmark} batas.
Penggabungan peringkat memakan puluhan mikrodetik, sedangkan memuat model
bahasa sesudah server dimulai ulang memakan hampir satu menit. Mempercepat
tahap penggabungan karena itu tidak akan terasa, dan kelompok yang memilihnya
sebagai sasaran perlu alasan lain di luar kecepatan.

\section{Benchmark Pembanding}

Angka hasil pengukuran tidak bermakna bila berdiri sendiri. Lama penggabungan
19,17 mikrodetik itu cepat atau lambat hanya dapat dijawab bila ada
pembandingnya. Karena itu setiap pengukuran wajib disandingkan dengan
sedikitnya satu \textit{benchmark}, dan \textit{benchmark} itu ditetapkan
sebelum pengukuran dimulai. \textit{Benchmark} yang dipilih sesudah hasilnya
keluar akan selalu menguntungkan penulisnya, dan penelaah mengenali pola itu.

Tiga jenis \textit{benchmark} dipakai pada proyek ini.
Tabel~\ref{tab:benchmark} memuat ketiganya menurut dimensi kualitasnya.

\begin{table}[htbp]
  \centering
  \caption{Tiga jenis \textit{benchmark} pembanding menurut dimensi yang diukur.}
  \label{tab:benchmark}
  \begin{tabularx}{\textwidth}{|l|X|X|X|}
    \hline
    \textbf{Dimensi} & \textbf{Dasar} & \textbf{Tandingan} & \textbf{Batas} \\
    \hline
    \textit{Maintainability} & Titik sentuh dan kopling Exim sebelum diubah & Modul lain pada Exim yang sudah memakai registri, misalnya \texttt{extractors} & Satu titik sentuh, yaitu penambahan yang hanya menyunting satu tempat \\
    \hline
    Performa & Lama tahap yang sama sebelum diubah, pada mesin dan beban yang sama & Pustaka mapan untuk pekerjaan yang sama, misalnya indeks vektor siap pakai & Biaya yang tetap dibayar meski pekerjaannya dihapus, misalnya lama membaca masukan saja \\
    \hline
    Ketahanan & Jumlah yang hilang pada susunan sebelum diubah & Jaminan yang dijanjikan rancangannya, misalnya menyambung dari nomor event terakhir & Nol kehilangan, yaitu seluruh event sampai meski koneksi diputus \\
    \hline
  \end{tabularx}
\end{table}

\textit{Benchmark} dasar wajib bagi seluruh judul, sebab tanpa keadaan awal
tidak ada yang dapat dibandingkan. Pengukurannya memakai skrip yang sama persis
dengan pengukuran sesudahnya, dan skrip itu tidak boleh ikut diubah di
antaranya.

\textit{Benchmark} tandingan menjaga agar klaimnya tidak berputar pada diri
sendiri. Susunan yang baru selalu lebih baik daripada susunan lama bila
pembandingnya hanya susunan lama itu. Tandingannya tidak harus dikalahkan,
sebab yang dilaporkan jaraknya, bukan kemenangannya.

\textit{Benchmark} batas menunjukkan ruang perbaikan yang sebenarnya tersisa.
Bab~\ref{bab:ruang} memakainya, yaitu antrean kosong melewatkan 1.200
\textit{request} dalam 15,5 ms, sehingga versi yang memakan 65,6 ms sudah
berada dekat batas itu. Angka seperti itu menghentikan pekerjaan yang sia-sia,
sebab memperbaiki bagian yang hanya menyumbang sepersekian dari keseluruhan
tidak akan terasa.

Ketiganya ditulis pada seksi Metode Penelitian, beserta alasan pemilihannya.
Seksi Hasil kemudian menyandingkan angka hasil pengukuran dengan ketiganya
pada tabel yang sama, sehingga pembaca tidak perlu membolak-balik halaman untuk
menilai besarannya.

\section{Struktur Makalah}

Makalah ditulis dalam bahasa Indonesia lebih dahulu, lalu diterjemahkan ke
bahasa Inggris sebelum dikirim, sebab SAUCIS hanya menerima naskah berbahasa
Inggris~\cite{saucis2026rules}. Kedua versi disimpan pada \textit{branch}
kelompok, sehingga riwayat penyuntingannya ikut terekam.

Urutan seksinya ditetapkan, dan Tabel~\ref{tab:struktur-makalah} memuat isi
tiap seksi. Seksi Metode Penelitian wajib menerangkan cara tiap \textit{research
question} dijawab, dan seksi Hasil wajib menjawabnya dengan urutan yang sama.

\begin{table}[htbp]
  \centering
  \caption{Struktur makalah beserta isi tiap seksinya.}
  \label{tab:struktur-makalah}
  \begin{tabularx}{\textwidth}{|l|X|}
    \hline
    \textbf{Seksi} & \textbf{Isinya} \\
    \hline
    Pendahuluan & Latar belakang dan motivasi, celah yang ditutup, \textit{research question}, dan tujuan \\
    \hline
    Landasan Teori & Gaya arsitektur yang dipakai beserta istilahnya, diambil dari bab yang relevan \\
    \hline
    Kajian Terkait & Pekerjaan orang lain pada masalah sejenis, beserta bedanya dengan pekerjaan ini \\
    \hline
    Metode Penelitian & Rancangan sebelum dan sesudah, ukuran yang dipakai, dan cara tiap RQ dijawab \\
    \hline
    Desain Solusi & Susunan yang diusulkan, beserta alasan tiap keputusannya \\
    \hline
    Evaluasi & Lingkungan percobaan, beban yang dipakai, jumlah putaran, dan cara menghitungnya \\
    \hline
    Hasil & Angka apa adanya, dalam tabel dan gambar, tanpa penafsiran \\
    \hline
    Pembahasan & Penafsiran angka itu, termasuk yang berlawanan dengan dugaan \\
    \hline
    \textit{Threats to Validity} & Ancaman terhadap kesahihan beserta cara menekannya \\
    \hline
    Kesimpulan dan \textit{Future Work} & Jawaban ringkas tiap RQ dan pekerjaan lanjutan yang masuk akal \\
    \hline
    Daftar Pustaka & Rujukan dengan gaya IEEE \\
    \hline
  \end{tabularx}
\end{table}

Seksi \textit{Threats to Validity} sering dilewati penulis pemula, padahal
seksi itulah yang paling diperiksa penelaah. Isinya menyebut hal yang dapat
membuat kesimpulan keliru, misalnya korpus tiruan yang tidak mewakili korpus
sebenarnya, atau pengukuran yang berjalan pada laptop yang juga menjalankan
pekerjaan lain~\cite{runeson2009case}.

\section{Aturan Template SAUCIS}

Naskah yang melanggar aturan penulisan tidak diteruskan ke penelaahan, sehingga
aturannya dibaca sebelum menulis, bukan sesudah.
Tabel~\ref{tab:aturan-saucis} memuat aturan yang paling sering terlewat.

\begin{table}[htbp]
  \centering
  \caption{Aturan penulisan SAUCIS yang wajib dipenuhi.}
  \label{tab:aturan-saucis}
  \begin{tabularx}{\textwidth}{|l|X|}
    \hline
    \textbf{Hal} & \textbf{Aturannya} \\
    \hline
    Bahasa & Bahasa Inggris, naskah berbahasa lain ditolak \\
    \hline
    \textit{File} & Template LaTeX, dan \textit{file} sumber LaTeX wajib diunggah \\
    \hline
    Kemiripan & Di bawah 20 persen, tidak termasuk daftar pustaka \\
    \hline
    Abstrak & Paling banyak 400 kata \\
    \hline
    Kata kunci & Tiga sampai lima kata kunci berbahasa Inggris, dipisah koma \\
    \hline
    Paragraf & Tanpa indentasi pada awal paragraf \\
    \hline
    Tabel & Berada di tengah, label dan keterangan di atas tabel \\
    \hline
    Gambar & Berada di tengah, label dan keterangan di bawah gambar \\
    \hline
    Persamaan & Berada di tengah, nomornya dalam kurung di tepi kanan \\
    \hline
    Sitasi & Gaya IEEE, bernomor dalam kurung siku \\
    \hline
    Panjang & Naskah eksperimental disarankan paling banyak 25 halaman dengan 7 gambar dan tabel \\
    \hline
    Unggahan & PDF beridentitas, PDF anonim, formulir hak cipta, dan \textit{file} sumber LaTeX \\
    \hline
  \end{tabularx}
\end{table}

\section{Trade-off yang Wajib Dilaporkan}

Setiap perbaikan arsitektur membawa harganya sendiri, dan makalah yang hanya
melaporkan sisi yang membaik terbaca sebagai promosi, bukan penelitian. Karena
itu tiap kelompok wajib menyebut \textit{trade-off} pekerjaannya pada seksi
Pembahasan, beserta angkanya, bukan sekadar menyebut bahwa harganya ada.

Harganya sering muncul pada dimensi yang berbeda dengan yang diperbaiki.
Perbaikan \textit{maintainability} biasanya dibayar dengan tambahan \textit{layer}
tak langsung, dan \textit{layer} itu menambah latensi. Perbaikan performa sering
dibayar dengan ketepatan, misalnya pencarian yang lebih cepat tetapi
hasilnya hampiran. Tabel~\ref{tab:tradeoff} memuat pasangan yang paling sering
muncul pada kedua puluh judul.

\begin{table}[htbp]
  \centering
  \caption{Pasangan perbaikan dan harganya yang wajib dilaporkan berangka.}
  \label{tab:tradeoff}
  \begin{tabularx}{\textwidth}{|X|X|X|}
    \hline
    \textbf{Yang diperbaiki} & \textbf{Harganya} & \textbf{Cara mengukur harganya} \\
    \hline
    Titik sentuh berkurang lewat registri & Satu \textit{layer} pencarian entri pada tiap pemanggilan & Latensi pemanggilan dengan dan tanpa registri, ratusan putaran \\
    \hline
    Pencarian lebih cepat lewat adapter hampiran & Hasilnya tidak lagi tepat & Irisan sepuluh teratas terhadap pencarian menyeluruh \\
    \hline
    Hasil naik lewat penambahan unit & Memori tiap unit tambahan dan biaya replikasi & Pemakaian memori per proses dan laju pada jumlah unit yang berbeda \\
    \hline
    Penulisan keluar dari jalur \textit{request} & Data yang sudah dijawab dapat lenyap & Jumlah yang hilang ketika satu unit dimatikan di tengah beban \\
    \hline
    \textit{File} dipecah agar tanggung jawabnya jelas & Pembaca baru menelusuri lebih banyak \textit{file} & Jumlah \textit{file} yang harus dibuka untuk satu perubahan \\
    \hline
    Inti bersama bagi beberapa turunan & Ketiganya terikat satu jadwal rilis & Jumlah pengujian yang ikut berjalan pada satu perubahan inti \\
    \hline
  \end{tabularx}
\end{table}

Satu bentuk pelaporan yang kuat adalah menyebut titik impasnya. Misalnya,
penambahan unit menguntungkan sampai empat unit lalu mendatar, atau
\textit{cache} menguntungkan ketika pengulangan pertanyaan melebihi sekian
persen. Titik impas seperti itu lebih berguna bagi pembaca daripada pernyataan
bahwa satu susunan lebih baik.

\section{Menaikkan Peluang Diterima}

Makalah yang hanya melaporkan satu sistem yang diperbaiki jarang lolos pada
jurnal bertingkat Q3. Sebabnya tiga. Objeknya tunggal, sehingga keberlakuan
umumnya lemah. Hasilnya dapat ditebak, sebab tidak ada penelaah yang meragukan
bahwa registri menurunkan titik sentuh. Artefaknya habis bersama naskahnya,
sehingga tidak ada yang dapat dipakai ulang orang lain.

Tiga cara menaikkannya, dan ketiganya terjangkau pada waktu satu semester.
Tabel~\ref{tab:peluang} merangkum ketiganya beserta tambahan pekerjaannya.

\begin{table}[htbp]
  \centering
  \caption{Tiga cara menaikkan peluang naskah diterima jurnal bertingkat.}
  \label{tab:peluang}
  \begin{tabularx}{\textwidth}{|l|X|X|}
    \hline
    \textbf{Cara} & \textbf{Isinya} & \textbf{Tambahan pekerjaannya} \\
    \hline
    Menambah objek & Ukuran yang sama dijalankan pada tiga sampai lima sistem sumber terbuka lain, bukan Exim saja & Paling murah, sebab skrip pembaca kode berjalan pada paket Python mana pun \\
    \hline
    Mengangkat ukurannya & Yang diusulkan bukan perbaikan satu sistem, melainkan satu ukuran beserta alatnya, lalu ditunjukkan bahwa ukuran itu meramalkan biaya sebenarnya & Menuntut pembuktian pada beberapa repositori sekaligus \\
    \hline
    Melaporkan \textit{trade-off} & Kedua sisi dilaporkan berangka, misalnya \textit{maintainability} membaik sekian sedangkan latensi naik sekian & Menuntut pengukuran performa, bukan pembacaan kode saja \\
    \hline
  \end{tabularx}
\end{table}

Perkiraan yang jujur perlu disampaikan di muka. Dengan cara pertama saja,
sasaran Q4 masuk akal. Sasaran Q3 menuntut cara pertama digabung salah satu
dari kedua cara lainnya. Tanpa ketiganya, sasaran yang wajar adalah jurnal
nasional terakreditasi atau konferensi, dan naskahnya tetap berguna sebagai
laporan studi kasus yang rapi~\cite{runeson2009case}.

\section{Kerja Kelompok dan Branch}

Proyek dikerjakan berkelompok. Tiap kelompok membuat satu \textit{branch} pada
repositori Exim, dan seluruh pekerjaannya tinggal di sana, yaitu perubahan
kode, skrip pengukuran, \textit{file} hasil pengukuran, serta naskah makalah kedua
bahasa. Riwayat penyuntingan makalah karena itu ikut terekam, sehingga
sumbangan tiap anggota terlihat.

Penamaan \textit{branch} mengikuti satu pola, yaitu nama berikut.

\begin{center}
  \texttt{capstone/<kelompok>-<judul-pendek>}
\end{center}

Isi direktorinya ditetapkan agar penilaian berjalan seragam, dan
Tabel~\ref{tab:isi-branch} memuat isinya.

\begin{table}[htbp]
  \centering
  \caption{Isi \textit{branch} tiap kelompok.}
  \label{tab:isi-branch}
  \begin{tabularx}{\textwidth}{|l|X|}
    \hline
    \textbf{Jalur} & \textbf{Isinya} \\
    \hline
    \texttt{app/} & Perubahan kode pada modul yang menjadi sasaran \\
    \hline
    \texttt{tests/} & Pengujian yang membuktikan perilakunya tidak berubah \\
    \hline
    \texttt{capstone/ukur/} & Skrip pengukuran beserta keluarannya apa adanya \\
    \hline
    \texttt{capstone/makalah/} & Naskah LaTeX berbahasa Indonesia dan terjemahannya \\
    \hline
    \texttt{capstone/README.md} & Judul, anggota kelompok, dan cara menjalankan pengukurannya \\
    \hline
  \end{tabularx}
\end{table}

Satu aturan tambahan berlaku pada \textit{commit}. Pengukuran keadaan awal
di-\textit{commit} sebelum perubahan kode dimulai, sehingga urutan waktunya
terbukti dari riwayat, bukan dari pengakuan. Keluaran pengukuran disimpan apa
adanya, termasuk kejanggalannya, lalu kejanggalan itu dijelaskan pada seksi
Pembahasan.

Pekerjaan dikumpulkan pada pertemuan terakhir, yaitu menurut jadwal Ujian
Akhir Semester. Yang dinilai isi \textit{branch} pada saat itu, sehingga
\textit{commit} yang datang sesudahnya tidak ikut dihitung. Karena riwayat
yang menjadi bukti, pekerjaan yang ditumpuk pada hari terakhir akan terlihat
dari tanggal \textit{commit}-nya.

\section{\textit{Summary}}

Bab ini memindahkan pelajaran empat belas bab sebelumnya ke satu sistem yang
benar-benar dipakai orang. Perbedaannya terasa sejak langkah pertama. Pada
contoh yang dibuat sendiri, gaya arsitekturnya dipilih lebih dahulu lalu
kodenya menyesuaikan. Pada sistem yang sudah berjalan, kodenya sudah ada lebih
dahulu, sehingga yang dapat dipilih hanya bagian mana yang diperbaiki dan
dengan teknik apa.

Pemilihan itu dimulai dari gejala yang terukur, bukan dari teknik yang sedang
diminati. Modul \texttt{config} yang diimpor 18 modul dari 23, dan modul
\texttt{server} yang mengimpor 12 modul, adalah gejala yang dapat dihitung
siapa pun yang menjalankan skripnya. Angka seperti itulah yang membuat latar
belakang sebuah makalah berpijak, sebab pembaca dapat memeriksanya sendiri.

Dua \textit{research question} menjadi tulang punggung naskah. Keduanya
menentukan ukuran yang dikumpulkan, menentukan isi Metode Penelitian, lalu
menentukan urutan seksi Hasil. Pertanyaan yang tidak menyebut ukurannya akan
melahirkan hasil yang tidak dapat dibantah sekaligus tidak dapat dipercaya,
dan naskah seperti itu berhenti di meja penelaah.

Kejujuran pengukuran dijaga tiga aturan. Pengulangan sedikitnya seratus kali,
pelaporan median beserta persentil ke-95 dan nilai terbesar, serta
perbandingan selisih dengan ragamnya. Ketiganya tampak merepotkan, namun
ketiganya pula yang membedakan klaim yang bertahan dari klaim yang runtuh pada
pertanyaan pertama penelaah.

Arsitektur yang baik, menurut ukuran bab ini, adalah arsitektur yang
perbaikannya dapat ditunjukkan dengan angka, dan angka itu dapat dihitung ulang
orang lain. Terdengar rapi saja belum cukup.

